\chapter{Anwendung des Style-Pakets} \index{Anwendung}
Nachfolgend ein kurzer Leitfaden f\"{u}r die Anwendung des
vorliegenden Style-Pakets f\"{u}r LaTeX-Dokumente f\"{u}r Arbeiten in der
Abteilung Informationssysteme

\begin{itemize}
\item Autor: Markus Schmees
\item Anpassungen f\"{u}r studentische Arbeiten: Marco Grawunder und Richard Hackelbusch
\item Version: 1.3.2
\item Stand: 24.04.2009
\end{itemize}

\section{Einbinden des Pakets}\index{Einbinden}\index{Paket}
Die einfachste M\"{o}glichkeit besteht darin, das vorgegebene Dokument
mit dem Titel

\texttt{BEARBEITER.tex}

zu verwenden und um eigene
Inhalte zu erg\"{a}nzen. BEARBEITER sollte durch die Art der Arbeit und den Namen des Autors ersetzt werden. Alternativ dazu kann der Autor im eigenen
\LaTeX-Hauptdokument, das die Dokumentenklasse \texttt{book}, eine
\texttt{11pt}-Schrift sowie die Blattgr\"{o}{\ss}e \texttt{a4paper}
verwenden muss, zus\"{a}tzlich am Ende der Paketdeklarationen das
Abteilungsstylepaket durch folgenden Befehl einbinden

\begin{verbatim}
\usepackage{abtisstud/abtisstud}
\end{verbatim}

Hierbei aber darauf achten, dass das Paket als \textbf{letztes
Paket} eingebunden wird, da es diverse \LaTeX-Befehle
\"{u}berschreibt, die von anderen Paketen ebenfalls \"{u}berschrieben
werden k\"{o}nnen. Das resultierende Hauptdokument sollte dann die in
der folgenden Abb.~\ref{stylepaket} dargestellte Struktur besitzen

\begin{figure}[ht]
\centering
\begin{Verbatim}[label=dissertation.tex,numberblanklines=false,fontsize=\scriptsize,numbers=left,frame=single]
\documentclass[11pt,a4paper]{book} % Basisdokumentenklasse
\usepackage[latin1]{inputenc}      % F\"{u}r Umlaute
\usepackage[T1]{fontenc}           % und {\ss}
\usepackage[ngerman]{babel}        % Deutsche Standardbezeichner und Trennung
\usepackage{makeidx}               % Indexpaket
\usepackage{times}                 % Sch\"{o}nere Schriften
\usepackage{longtable}             % Lange Tabelle f\"{u}r Abk\"{u}rzungsverzeichnis
\usepackage{graphics}              % EPS-Grafiken einbinden
\usepackage{abtisstud/abtisstud}   % Das Stylepaket am Ende einbinden!

% Das eigentliche Dokument
\begin{document}
... Inhalt ...
\end{document}
\end{Verbatim}
\caption{Einbinden des Style-Pakets}\label{stylepaket}
\end{figure}


\section{\"{A}nderungen gegen\"{u}ber \texttt{book}}\index{Anderungen@\"{A}nderungen}
Dem Erweiterungspaket zugrunde liegt der Dokumentenstil
\texttt{book}. Dieser wurde mit Hilfe des vorliegenden Pakets an
Besonderheiten der IS-Formatvorlage angepasst. Dabei wurden
folgende Dinge ver\"{a}ndert:

\begin{itemize}
\item Abschnittnummerierung auf Ebene 3 gesetzt (bis 1.1.1.1)
\item Auff\"{u}hrung im Inhaltsverzeichnis bis Ebene 1 (bis 1.1)
\item R\"{a}nder, Header, Textweite und Texth\"{o}he angepasst
\item Linken Einzug um 0,37cm bei neuem Absatz eingef\"{u}gt
\item Abstand der Abs\"{a}tze auf 6pt eingestellt
\item Seitenzahlen und Abschnitte im Header dargestellt
\item Abstand der \"{U}berschriften zum Text eingestellt
\item Farbe und Schrift der \"{U}berschriften festgelegt
\item Verzeichnisdeklarationen f\"{u}r korrekte Headeranzeige
\"{u}berschrieben
\item Abstand der Abbildungen zu den Unterschriften angepasst
\item Schriftart und Style f\"{u}r Abbildungsunterschriften
eingestellt
\item "`Offizielle"' Titelseite integriert
\item Seitenzahltypen (r\"{o}misch und arabisch) kombiniert
\item Fu{\ss}noten linksb\"{u}ndig und mit Abstand zwischen Nummer und
Text
\item S\"{a}mtliche Listen linksb\"{u}ndig
\item Geringere Einr\"{u}ckung der Unterstichpunkte im Index
\item Umwandlung der \"{U}berschriften f\"{u}r Verzeichnisse in
gleichwertige \"{U}berschriften ohne "`-ver\-zeich\-nis"'
\end{itemize}


\section{Einbinden von Abbildungen}\index{Abbildungen}
\index{Abbildungen}
\index{Einbinden!Abbildungen}
Wie bei der normalen Verwendung von \LaTeX{} k\"{o}nnen auch hier \"{u}ber
die entsprechenden Zusatzpakete Grafiken eingebunden werden. Immer
darauf achten, dass diese als Vektorgrafiken vorliegen, damit sie
problemlos skaliert und damit auch ohne Qualit\"{a}tsverlust gedruckt
werden k\"{o}nnen. Nachfolgend zeigt Abb.~\ref{abbildung} am Beispiel
des Uni-Logos, wie eine Abbildung eingebunden und angezeigt werden
kann (vlg. dazu auch Abschnitt \ref{Schriften}).

\begin{figure}[ht]
  \centering
  \includegraphics[scale=1]{style/unilogo.pdf}
  \caption{Eine sch\"{o}ne Abbildung}\label{abbildung}
\end{figure}


\section{Erstellen und Anpassen des Indexes}\index{Indexerstellung}
Um das grundlegende Indexfile zu erzeugen, muss der Befehl
\texttt{makeindex} innerhalb des Haupt-\LaTeX-Dokuments aufgef\"{u}hrt sein.
Dieser erzeugt aus s\"{a}mtlichen \texttt{index}-Befehlen die Datei

\texttt{dissertation.idx}.

Um diese Datei ebenfalls passend zu
formatieren, wurde ein entsprechendes Stylefile namens

\texttt{abtisstud.ist}

erstellt. Die Anwendung dieses Styles auf den
Index erfolgt mit Hilfe des folgenden Befehls

\begin{verbatim}
makeindex -s abtisstud/abtisstud.ist -o content/index.tex \
    dissertation.idx
\end{verbatim}

Dieser Befehl sorgt daf\"{u}r, dass die Indexdaten entsprechend
formatiert und in die Datei \texttt{index.tex} in das Verzeichnis
\texttt{content} abgelegt werden. Von dort kann der neue Index
dann problemlos per input eingelesen werden. Also bitte nicht den
traditionell erzeugten Index \"{u}ber den Befehl
\texttt{printindex} einbinden! Der Auszug aus dem
\LaTeX-Hauptdokument zur Erzeugung und Einbindung des Indexes sieht
folgenderma{\ss}en aus

\begin{figure}[ht]
\centering
\begin{Verbatim}[label=dissertation.tex,numberblanklines=false,fontsize=\scriptsize,numbers=left,frame=single]
...

% Indexdatei erzeugen
\makeindex

% Das eigentliche Dokument
\begin{document}

...

\index{Irgendwas}                     % Einstufiger Indexeintrag
\index{Irgendwas!Ganzanderes}         % Zweistufiger Indexeintrag

...

% Index als letztes Verzeichnis
\cleardoublepage                      % Auf ungerader Seite beginnen
\addcontentsline{toc}{chapter}{Index} % Index im Inhaltsverzeichnis anzeigen
\begin{theindex}
{\indexfrontskip\Large\sffamily\bfseries\hfill A\hfill}\nopagebreak
 
  \item Abbildungen\dotfill 2
  \item \"Anderungen\dotfill 2
  \item Anwendung\dotfill 1

  \indexspace
{\indexfrontskip\Large\sffamily\bfseries\hfill E\hfill}\nopagebreak
 
  \item Einbinden\dotfill 1
    \subitem Abbildungen\dotfill 2

  \indexspace
{\indexfrontskip\Large\sffamily\bfseries\hfill F\hfill}\nopagebreak
 
  \item Festlegungen\dotfill 7

  \indexspace
{\indexfrontskip\Large\sffamily\bfseries\hfill I\hfill}\nopagebreak
 
  \item Indexerstellung\dotfill 3
  \item Internationalisierung\dotfill 4

  \indexspace
{\indexfrontskip\Large\sffamily\bfseries\hfill L\hfill}\nopagebreak
 
  \item Literatur\dotfill 7
  \item Literaturreferenzen\dotfill 7

  \indexspace
{\indexfrontskip\Large\sffamily\bfseries\hfill P\hfill}\nopagebreak
 
  \item Paket\dotfill 1

  \indexspace
{\indexfrontskip\Large\sffamily\bfseries\hfill S\hfill}\nopagebreak
 
  \item Schriften\dotfill 5

  \indexspace
{\indexfrontskip\Large\sffamily\bfseries\hfill T\hfill}\nopagebreak
 
  \item Titelseite\dotfill 8

  \indexspace
{\indexfrontskip\Large\sffamily\bfseries\hfill V\hfill}\nopagebreak
 
  \item Verzeichnisreihenfolge\dotfill 7

  \indexspace
{\indexfrontskip\Large\sffamily\bfseries\hfill Z\hfill}\nopagebreak
 
  \item Zitate\dotfill 7

\end{theindex}
             % Formatierten Index einbinden
\end{document}
\end{Verbatim}
\caption{Index erstellen, formatieren und integrieren}\label{indexintegration}
\end{figure}


\section{Internationalisierung}\index{Internationalisierung}
Durch Einbinden entsprechender Sprachpakete wird \LaTeX{} an
Besonderheiten bestimmter Sprachr\"{a}ume, z.B. deren Trennung oder
Verzeichnisbezeichnung angepasst. Durch die individuelle
Gestaltung bestimmter Verzeichnisnamen wurden die zugeh\"{o}rigen
Befehle aber \"{u}berschrieben. Das bedeutet nun, dass der vorliegende
Stil zwar problemlos f\"{u}r deutsche Arbeiten anwendbar ist, aber
diese Verzeichnisbezeichnung f\"{u}r Dissertationen in fremder Sprache
angepasst werden m\"{u}ssen. Diese betreffen also insbesondere die
folgenden Paket- und Verzeichnisdeklarationen in den Zeilen 3, 11,
14, 18, 19, 23, 25, 29 und 30.

\begin{figure}[htp]
\centering
\begin{Verbatim}[label=dissertation.tex,numberblanklines=false,fontsize=\scriptsize,numbers=left,frame=single]
\documentclass[11pt,a4paper]{book}

...

\usepackage[ngerman]{babel}

...

\usepackage{abtisstud/abtisstud}

...

% Das eigentliche Dokument
\begin{document}

...

% Verzeichnisse am Ende, erst das Glossar
\addonchapter{Glossar} %
Nachfolgend sind noch einmal wesentliche Begriffe dieser Arbeit
zusammengefasst und erl�utert. Eine ausf�hrliche Erkl�rung findet
sich jeweils in den einf�hrenden Abschnitten sowie der jeweils
darin angegebenen Literatur. Das im Folgenden im Rahmen der
Erl�uterung verwendete Symbol \this bezieht sich jeweils auf den
im Einzelnen vorgestellten Begriff, das Symbol \siehe{} verweist
auf einen ebenfalls innerhalb dieses Glossars erkl�rten Begriff.

\begin{description}
\item[Auktion] Eine \this ist das im \siehe{E-Commerce} am
H�ufigsten eingesetzte Verfahren zur dynamischen
\siehe{Preisfindung}. Interessenten k�nnen dabei durch Abgabe von
Geboten Preis, Dauer und Gewinner beeinflussen. Bei einer offenen
\this sind Bieter, H�he der Gebote und der aktuelle Preis f�r
alle Teilnehmer sichtbar, bei der geschlossenen (sealed) \this
erfolgt nur eine interne Benachrichtigung. Die bekanntesten Typen
sind die traditionelle \siehe{Versteigerung} sowie die
\siehe{holl�ndische}, \siehe{umgekehrte} und \siehe{verdeckte} \this.

\item[Behaviorismus] Der \this ist eine \siehe{Lerntheorie}, die
davon ausgeht, dass Wissen als Struktur unabh�ngig vom
\siehe{Lernenden} existiert und dass sein Verhalten operant
konditioniert ist, d.h. dass es als Konsequenz aus anderen
Verhaltensweisen resultiert. Erfolgt eine positive Reaktion,
beh�lt der \siehe{Lernende} neu erlerntes Verhalten bei, negative
Reaktionen f�hren zu einer Verminderung dieses Verhaltens. Der
\siehe{Lehrende} bestimmt dabei das zu erlernende Wissen und
ist f�r die Steuerung des \siehe{Lernprozesses} zust�ndig.
\end{description}


% Dann die Abk\"{u}rzungen
\addonchapter{Abk\"{u}rzungen}
\begin{longtable}[ht]{ll}
ADDIE & Analysis, Design, Development, Implementation, Evaluation\\  %E-Learning
ADEPT & Advanced Decision Environment for Process Tasks\\  %Konzept
ADL & Advanced Distributed Learning Initiative\\  %E-Learning
AGB & Allgemeine Gesch�ftsbedingungen\\  %Konzept
AGOF & Arbeitsgemeinschaft Online Forschung\\  %Einleitung
AICC & Aviation Industry Computer Based Training Committee\\  %E-Learning
API & Application Programming Interface\\  %TEL, Konzept
%ARIADNE & Alliance of Remote Instructional Authoring and Distribution Networks for Europe\\  %E-Learning
ARIS & Architektur integrierter Informationssysteme\\  %Konzept
ASC & Accredited Standards Committee\\  %EC-Grundlagen
ASP & Application Service Providing\\  %Integration
ASTD & American Society for Training and Development\\  %E-Learning
AXIS & Apache eXtensible Interaction System\\  %Implementierung
B2B & Business-to-Business\\  %Konzept, Glossar
B2C & Business-to-Consumer\\  %Konzept, Glossar
BDSG & Bundesdatenschutzgesetz\\  %EC-Grundlagen
BGB & B�rgerliches Gesetzbuch\\  %Konzept
BITKOM & Bundesverband Informationswirtschaft, Telekommunikation und neue Medien\\  % Grundlagen E-Commerce
BMBF & Bundesministerium f�r Bildung und Forschung\\  %Einleitung
BME & Bundesverband Materialwirtschaft, Einkauf und Logistik\\  %EC-Grundlagen
BPEL4WS & Business Process Execution Language for Web Services\\  %Implementierung
BSCW & Basic Support for Cooperative Work\\  %E-Learning
BSI & Bundesamt f�r Sicherheit in der Informationstechnik\\ %Literatur
CAL & Computer Aided/Assisted Learning\\  %E-Learning
CBT & Computer Based Training\\  %E-Learning, Vorgaben
CC & Creative Commons\\  %Konzept
CD & Compact Disc\\  %E-Learning
CELab & Labor f�r Content Engineering\\  %TEL
CMI & Computer Managed Instruction\\  %E-Learning
CMS & Content Management System\\  %TEL, Konzept
CORBA & Common Object Request Broker Architecture\\  %Konzept
CPU & Central Processing Unit\\  %Konzept
CSCL & Computer Supported Collaborative Learning\\  %E-Learning
CSCW & Computer Supported Cooperative Work\\  %TEL
CSS & Customer Support Services\\  %Konzept
CRM & Customer Relationship Management\\  %Konzept
CUL & Computerunterst�tztes Lernen\\  %E-Learning
DBMS & Datenbankmanagementsystem\\  %Konzept
DCOM & Distributed Component Object Model\\  %Konzept
DFN & Deutsches Forschungsnetz\\  %E-Learning
DIN & Deutsches Institut f�r Normung\\  %Lernen, E-Learning, TEL
DREL & Digital Rights Expression Language\\  %Konzept
DRM & Digital Rights Management\\  %Konzept
DVD & Digital Video Disc\\  %E-Learning
E2B & Education-to-Business\\  %Integration
E2C & Education-to-Consumer\\  %Integration
E2E & Education-to-Education\\  %Integration
EAI & Enterprise Application Integration\\  %Konzept
EAN & European Article Numbering\\  %Grundlagen-EC-Standards
EBPP & Electronic Bill Presentment and Payment\\  %Konzept
ebXML & Electronic Business Extensible Markup Language\\ %EC-Grundlagen
ECA & Event-Condition-Action\\  %Vorgaben
EC & Electronic Cash\\  %EC-Grundlagen, Konzept
ECC & E-Learning Courseware Certification\\  %E-Learning
EDI & Electronic Data Interchange\\  %EC-Grundlagen
EDIFACT & EDI for Administration, Commerce and Transport\\  %Konzept
EFQM & European Foundation for Quality Management\\  %TEL
EGBGB & Einf�hrungsgesetz zum B�rgerlichen Gesetzbuch\\  %Konzept
EITO & European Information Technology Observatory\\  % Grundlagen E-Commerce
ELAN & E-Learning Academic Network Niedersachsen\\  %Einleitung, Vorgaben
\end{longtable}


% Weiter mit Abbildungen
\cleardoublepage
\addcontentsline{toc}{chapter}{Abbildungen}
\renewcommand{\listfigurename}{Abbildungen}
\listoffigures

% Schlie{\ss}lich Literatur
\cleardoublepage
\addcontentsline{toc}{chapter}{Literatur}
\bibliographystyle{alphadin}
\renewcommand{\bibname}{Literatur}
\bibliography{content/bibliographie}

% Und ganz am Ende der Index
\cleardoublepage
\addcontentsline{toc}{chapter}{Index}
\begin{theindex}
{\indexfrontskip\Large\sffamily\bfseries\hfill A\hfill}\nopagebreak
 
  \item Abbildungen\dotfill 2
  \item \"Anderungen\dotfill 2
  \item Anwendung\dotfill 1

  \indexspace
{\indexfrontskip\Large\sffamily\bfseries\hfill E\hfill}\nopagebreak
 
  \item Einbinden\dotfill 1
    \subitem Abbildungen\dotfill 2

  \indexspace
{\indexfrontskip\Large\sffamily\bfseries\hfill F\hfill}\nopagebreak
 
  \item Festlegungen\dotfill 7

  \indexspace
{\indexfrontskip\Large\sffamily\bfseries\hfill I\hfill}\nopagebreak
 
  \item Indexerstellung\dotfill 3
  \item Internationalisierung\dotfill 4

  \indexspace
{\indexfrontskip\Large\sffamily\bfseries\hfill L\hfill}\nopagebreak
 
  \item Literatur\dotfill 7
  \item Literaturreferenzen\dotfill 7

  \indexspace
{\indexfrontskip\Large\sffamily\bfseries\hfill P\hfill}\nopagebreak
 
  \item Paket\dotfill 1

  \indexspace
{\indexfrontskip\Large\sffamily\bfseries\hfill S\hfill}\nopagebreak
 
  \item Schriften\dotfill 5

  \indexspace
{\indexfrontskip\Large\sffamily\bfseries\hfill T\hfill}\nopagebreak
 
  \item Titelseite\dotfill 8

  \indexspace
{\indexfrontskip\Large\sffamily\bfseries\hfill V\hfill}\nopagebreak
 
  \item Verzeichnisreihenfolge\dotfill 7

  \indexspace
{\indexfrontskip\Large\sffamily\bfseries\hfill Z\hfill}\nopagebreak
 
  \item Zitate\dotfill 7

\end{theindex}


% Und schlie{\ss}lich noch die Versicherung f\"{u}r das Pr\"{u}fungsamt
\cleardoublepage
\versicherung{Oldenburg}

\end{document}
\end{Verbatim}
\caption{Anzupassende Pakete bzw. Verzeichnisbezeichner}\label{verzeichnisanpassung}
\end{figure}




\section{Einbetten s\"{a}mtlicher Schriften}\index{Schriften}\label{Schriften}
Bitte unbedingt folgende Dokumenteneinstellung beibehalten

\begin{verbatim}
\documentclass[11pt,a4paper]{book}
\end{verbatim}

Grunds\"{a}tzlich bietet es sich an, die Arbeit mit \texttt{pdflatex} direkt als PDF zu erzeugen. Bei der Verwendung von \texttt{pdflatex} k\"{o}nnen als Bildformate PDF, PNG oder JPG (Achtung: Die beiden letzteren sind nicht im Vektorformat und damit nicht gut f\"{u}r den Ausdruck geeignet.) f\"{u}r Bilder verwendet werden.

Es sollten aber s\"{a}mtliche verwendete Schriftarten in das
resultierende PDF-Do\-kument eingebettet werden. \texttt{pdflatex}
und \texttt{ps2pdf} sparen unter Umst\"{a}nden einige Schriften aus, sollten also nur mit Bedacht verwendet werden.

Eine M\"{o}glichkeit, dennoch s\"{a}mtliche Schriftarten in das
resultierende Dokument einzubetten, bietet eine Kombination aus
\texttt{dvips} und \texttt{ghostscript} an. Zun\"{a}chst das
resultierende DVI-Dokument \texttt{dissertation.dvi} in ein
passendes PS-Dokument umwandeln mit dem folgenden Befehl

\begin{verbatim}
dvips -P pdf -t A4 -z -Pdownload35 dissertation.dvi
\end{verbatim}

Die resultierende Datei \texttt{dissertation.ps} darf daraufhin
NICHT mit
\texttt{ps2pdf} sondern muss -- wie der folgende Befehl zeigt --
mit \texttt{ghostscript} in ein PDF umgewandelt werden

\begin{verbatim}
gs -dSAFER -dNOPAUSE -dBATCH -sDEVICE=pdfwrite \

   -sPAPERSIZE=a4 -dPDFSETTINGS=/printer \

   -dCompatibilityLevel=1.3 -dMaxSubsetPct=100 \

   -dSubsetFonts=true -dEmbedAllFonts=true \

   -sOutputFile=dissertation.pdf dissertation.ps
\end{verbatim}

Das hieraus resultierende PDF-Dokument \texttt{dissertation.pdf}
enth\"{a}lt schlie{\ss}lich sowohl das gew\"{u}nschte Format als auch
s\"{a}mtliche Schriften. Bei dem Umweg \"{u}ber EPS ist allerdings zu beachten, dass Vektorgraphiken nicht als PDF sondern als EPS-Dateien eingebunden werden m\"{u}ssen.

\chapter{Festlegungen der Abteilung}\index{Festlegungen}
\"{U}ber das reine Design hinaus sind nachfolgend Festlegungen zur
Struktur einer Arbeit (insbesondere Studentischer Arbeiten) der
Abteilung Informationssysteme  aufgef\"{u}hrt.



\section{Verzeichnisreihenfolge}\index{Verzeichnisreihenfolge}
Die Reihenfolge der Verzeichnisse soll dem vorliegenden Dokument
folgen, die Verzeichnisse dazu wie folgt angeordnet sein. Im
Inhaltsverzeichnis, das neuerdings nur noch "`Inhalt"' hei{\ss}t,
werden nur die nachfolgend fett dargestellten Punkte aufgef\"{u}hrt.

\begin{enumerate}
\item Zusammenfassung
\item Abstract
\item Inhalt (vormals Inhaltsverzeichnis)
\item \textbf{Der eigentliche Inhalt} mit laufenden Kapitelnummern
\item \textbf{Glossar} und Folgende jeweils ohne Kapitelnummer
\item \textbf{Abk\"{u}rzungen}
\item \textbf{Abbildungen}
\item \textbf{Literatur}
\item \textbf{Index}
\end{enumerate}


Hierbei wurde auf die traditionelle Verwendung der Begriffe
Inhaltsverzeichnis, Literaturverzeichnis und Abbildungsverzeichnis
verzichtet und stattdessen nur noch Inhalt, Literatur und
Abbildungen verwendet.


\section{Zitate und Literaturreferenzen}\index{Zitate}\index{Literaturreferenzen}\index{Literatur}
Im eigentlichen Text kann man verschiedene Zitate verwenden. Das
machen auch z.B. \cite{abel03}, \cite{appelrath06}, \cite{back01}
oder \cite{bakos91}. Den jeweiligen Zitaten soll dazu der
Bibliographiestil \texttt{alphadin} (mit kleinem "`a"') zugrunde
liegen. Dieser repr\"{a}sentiert einen genormten Zitierstil und wird
durch folgende Zeile im \LaTeX-Dokument eingestellt

\begin{verbatim}
\bibliographystyle{alphadin}
\end{verbatim}

Wer trotz des verk\"{u}rzten Zitierstils in ausf\"{u}hrlicher Weise auf
die jeweiligen Autoren hinweisen m\"{o}chte, hat die M\"{o}glichkeit, dies
z.B. folgenderma{\ss}en zu machen: "`wie auch von Appelrath, Boles,
Kleinefeld und andere in
\cite{appelrath06} beschrieben\ldots"'

Bei den Literaturangaben wird k\"{u}nftig auf die ISBN- bzw.
ISSN-Nummern verzichtet, da bei Dissertationen ohnehin davon
auszugehen ist, dass begutachtete bzw. verlegte Literatur
verwendet wird.


\section{Die offizielle Titelseite}\index{Titelseite}
Der \texttt{maketitle}-Befehl wurde \"{u}berschrieben. F\"{u}r die neue Titelseite
sind daher folgende Einstellungen zu machen

\begin{figure}[ht]
\centering
\begin{Verbatim}[label=dissertation.tex,numberblanklines=false,fontsize=\scriptsize,numbers=left,frame=single]
% Voreinstellungen f\"{u}r offizielle Titelseite
\title{Hier den Titel einf\"{u}gen}                   % Titel der Dissertation
\author{Dipl.-Inform. Vorname Nachname}           % Name des Autors
\arbeitstyp{Diplomarbeit}                         % oder Bachelorarbeit, Masterarbeit ...
\erstgutachter{Prof. Dr. Erster Gutachter}        % Name des Erstgutachters
\zweitgutachter{Prof. Dr.-Ing. Zweiter Gutachter} % Name des Zweitgutachters

% Einbinden/Anzeigen der Titelseite innerhalb von \begin/end{document}
\maketitle
\end{Verbatim}
\caption{Einstellungen f\"{u}r offizielle Titelseite}\label{titelseite}
\end{figure}
